\section*{Abstract}

Particle physics is at a point of a very statistical and gradual unrevealing of the future of science. Most of the processes that are of interest in our discipline occur within time and length scales that are practically impossible to record or observe. The strong and electroweak interactions that we are trying to understand can only be probed by the rigorous analysis of the final and initial states of the particles involved in these interactions. The strategy scientists decided to employ in order to study these fundamental interactions has evolved ever so slightly since the 50's, when CERN was founded. The true evolution occurred in our detector and computational capacity. Today, in the Large Hadron Collider (LHC), we are colliding protons with 13 TeV center-of-mass energies at \( 4 \times 10^{8} \) times per second. A high center-of-mass energy and a frequent collision rate unlocks the possibility of analyzing rarer processes or those that are suppressed at lower energies. This increased capacity comes with a caviat; the perfect physics breeding ground we have at our collision point results in huge amounts of physical processes to occur per collision and accordingly, massive amount of combination data (including background and unresolvable events) is registered for each collision event.\\

Our purpose built detectors manage to record the final states of the fundamental particle interactions very well; however, the flipside of the coin, the massive data that we collect needs to be analyzed with many layers of similarly purposefully designed hardware systems, software systems and algorithms. Usually, we are interested in figuring out the details of a specific decay process. The detector measurement for any extremely rare hard scattering process in LHC is buried under statistical uncertainty of every other process, such as QCD multi-jet events, therefore we need to come up with a multi-layered strategy to systematically select only the signal for the process we want to study. In this laboratory course, we are going to conduct a comprehensive data analysis where we will look for a theoretical extension to the Standard Model "$Z'$" particle decaying into a "$t\overline{t}$" pair, utilizing the data from the ATLAS detector at LHC.